\begin{frame}{$\beta$-VAE: 두 항의 균형점을 ``학습률처럼''}
  복원 항과 정규화 항의 계수를 1:1로 고정한 것은 하나의 선택일 뿐.
  $\beta$-VAE (H. Kingma, Salimans, Welling, 2016)는 정규화 항에
  계수 $\beta$를 붙여 $\mathcal{L} = \mathbb{E}[\log p(x|z)]
  - \beta\, D_{KL}$로 일반화:
  \vskip0.6em
  \small
  \begin{tabular}{@{}lrrr@{}}
    \toprule
    $\beta$ & 복원 손실 & KL & 생성 샘플 표준편차 \\
    \midrule
    0   & 0.016 & 19.1 & 0.58 \\
    0.1 & 0.068 & 5.61 & 0.91 \\
    0.5 & 0.232 & 1.70 & 1.67 \\
    1.0 (기본) & 0.378 & 1.22 & 2.05 \\
    2.0 & 0.546 & 0.90 & 2.42 \\
    \bottomrule
  \end{tabular}
  \vskip0.8em
  \begin{itemize}
    \item $\beta = 0$: 정규화 항 \textbf{완전히 끔} → 일반 오토인코더로
      복귀. 복원 손실 0.016(아주 작음)이지만 KL 19.1(잠재 공간이
      표준정규와 거의 무관), 생성 표준편차 0.58 — \textbf{데이터(표준편차
      4.05)보다 훨씬 좁음}. 잠재 공간이 ``여기저기 흩어져 있어 아무 데서
      샘플링해도 안 된다''는 원래 문제 그대로.
    \item $\beta = 2$: 정규화 항 강화 → KL 0.90으로 작아지지만 복원
      손실 0.55로 올라가고, 생성 표준편차 2.42 — \textbf{정보를 더
      압축}하는 대가로 복원 품질이 약간 떨어짐.
  \end{itemize}
  \vskip0.4em
  {\footnotesize $\beta$는 ``복원 품질 vs 잠재 공간의 매끄러움(정규성)''을
  \textbf{학습률처럼 튜닝하는 하이퍼파라미터}. (``VAE가 왜 흐릿한지''의
  답변도 결국 이 $\beta$의 균형 문제에서 출발.)}
\end{frame}
